\section{Conclusion}
Neural Bellman Operators construct attained continuations and feasible dynamic policies with an original-optimum error account. Original-owner compilation and reference-centered recursive-risk qualification preserve the economic object while exposing the work required to evaluate it. The continuous investment construction supplies the necessary curvature, supports and full-law inequalities from economic primitives.

The additional action-aligned result identifies the quantities that control a proposal's decision gap: the width of an original-objective enclosure, the directional support mismatch and the constrained stationarity residual. Prospective selection can then target rejection or complete work rather than infer either from a smaller global residual span. The exact purchase mechanism demonstrates a strictly positive economic value of an oblique continuation relative to the best specified additive and quadratic classes, while equal conventional geometry attains equal value. Nonuniform-law and explicitly quantified randomized information results place the precision account on a wider mathematical basis.

The investment tournament remains an essential discipline on these constructions. Its quadratic least-squares winner, increased aggregate rejections under residual-oriented fitting, and absence of certified neural adoption over that winner are preserved. The mathematical mechanisms and that empirical comparison answer different questions. Together they identify what an NBO must construct, what its certificate must cover and what evidence is needed to establish a useful decision at a competitive complete cost.
